\chapter{Concluzii}
\label{ch:concluzii}

Întrebarea formulată în Capitolul~\ref{ch:introducere} a fost în ce măsură modul în care un model reprezintă structura latentă a datelor \enquote{normale} determină tipurile de anomalii pe care le poate detecta. Rezultatele din Capitolul~\ref{ch:results}, confirmate și nuanțate în Capitolul~\ref{ch:explainability}, dau un răspuns afirmativ: aceeași paradigmă produce rezultate diferite pe seturi de date diferite, nu pentru că modelul ar fi mai bun sau mai slab în sine, ci pentru că presupunerea lui despre normalitate se potrivește sau nu structurii reale a anomaliei căutate.

Cât despre cele două ipoteze formulate în \S\ref{sec:expected-behaviour}, ele se confirmă, dar nu necondiționat. Prima ipoteză -- că structura temporală inerentă avantajează seturile de date cu adevărat secvențiale -- se confirmă clar pe Yahoo A3/A4, unde paradigma predictivă domină, dar se inversează pe Setul de Date A (unde secvența este artificială) și este depășită de caracteristicile de context la CIC-IDS2017, unde reconstrucția câștigă în ciuda structurii temporale disponibile.\
A doua ipoteză -- că paradigme complementare, combinate într-un ansamblu, îmbunătățesc performanța -- este confirmată de studiul de caz din \S\ref{sec:explain-disjoint-regime}: AE+Context și Temporal IForest exploatează, demonstrabil prin SHAP, evidențe disjuncte, iar votarea ponderată extinde acoperirea pe clase de atac pe care niciun model singur nu le detectează bine. Câștigul nu este însă gratuit -- vine cu un cost pe AP-ul agregat, dominat de clasele majoritare -- ceea ce nuanțează, mai degrabă decât infirmă, ipoteza inițială.

Limitările acestei lucrări deschid două direcții clare de continuare: testarea sistematică a celorlalte presupuneri ridicate în discuțiile Capitolului~\ref{ch:results} -- de exemplu de ce ferestrele temporale dăunează performanței la Setul de Date A și C, sau de ce clase precum FTP-Patator, SSH-Patator, Bot rămân nedetectabile indiferent de setul de caracteristici; și extinderea studiului de ansamblu din \S\ref{sec:explain-ensemble-build} la Setul de Date A și B, cu reguli de combinare mai sofisticate decât votarea ponderată și regula maximului.

\enlargethispage{\baselineskip}